\documentclass[11pt,a4paper]{article}

% ---------- Encodage & langue ----------
\usepackage[utf8]{inputenc}
\usepackage[T1]{fontenc}
\usepackage[french]{babel}
\usepackage{lmodern}

% ---------- Mise en page ----------
\usepackage[margin=2.3cm]{geometry}
\usepackage{setspace}
\onehalfspacing
\usepackage{parskip}

% ---------- Couleurs & graphismes ----------
\usepackage{xcolor}
\definecolor{primary}{HTML}{1F6FEB}
\definecolor{dark}{HTML}{0D1117}
\definecolor{lightgray}{HTML}{F0F3F7}
\definecolor{success}{HTML}{2EA043}
\definecolor{warning}{HTML}{D29922}

% ---------- Tableaux ----------
\usepackage{array}
\usepackage{booktabs}
\usepackage{longtable}
\usepackage{tabularx}
\usepackage{colortbl}
\renewcommand{\arraystretch}{1.3}

% ---------- Titres de sections ----------
\usepackage{titlesec}
\titleformat{\section}
  {\color{primary}\Large\bfseries}{\thesection.}{0.6em}{}
\titleformat{\subsection}
  {\color{dark}\large\bfseries}{\thesubsection}{0.6em}{}
\titlespacing*{\section}{0pt}{1.4em}{0.7em}

% ---------- En-têtes / pieds de page ----------
\usepackage{fancyhdr}
\pagestyle{fancy}
\fancyhf{}
\lhead{\small\textcolor{primary}{Medical Pipeline}}
\rhead{\small Rapport de fin d'année}
\cfoot{\thepage}
\renewcommand{\headrulewidth}{0.4pt}

% ---------- Liens ----------
\usepackage{hyperref}
\hypersetup{
  colorlinks=true,
  linkcolor=primary,
  urlcolor=primary,
  pdftitle={Rapport de fin d'annee - Medical Pipeline},
  pdfauthor={ElKarizma}
}

% ---------- Icônes texte ----------
\newcommand{\ok}{\textcolor{success}{\textbf{OK}}}
\newcommand{\partiel}{\textcolor{warning}{\textbf{Partiel}}}

\begin{document}

% ================= PAGE DE TITRE =================
\begin{titlepage}
\centering
\vspace*{2cm}

{\color{primary}\rule{\linewidth}{2pt}}\\[0.6cm]
{\Huge\bfseries Rapport de Fin d'Année\par}
\vspace{0.5cm}
{\LARGE Projet \textit{Medical Pipeline}\par}
\vspace{0.3cm}
{\large Diagnostic IA à partir des symptômes\\ Scraping NHS \& Machine Learning\par}
{\color{primary}\rule{\linewidth}{2pt}}\\[1.5cm]

\begin{tabular}{rl}
\textbf{Projet} & Medical Pipeline -- Diagnostic IA NHS \\[4pt]
\textbf{Auteur} & ElKarizma \\[4pt]
\textbf{Période} & Année 2025 \\[4pt]
\textbf{Date du rapport} & 10 juin 2026 \\[4pt]
\textbf{Dépôt} & \texttt{ElKarizma/medical\_pipeline-master} \\[4pt]
\textbf{Stack} & Python, Streamlit, scikit-learn, BeautifulSoup, DVC \\
\end{tabular}

\vfill
{\small\itshape Outil à finalité éducative -- ne se substitue pas à un avis médical professionnel.\par}
\end{titlepage}

\tableofcontents
\newpage

% ================= 1. RESUME =================
\section{Résumé exécutif}

Le projet \textbf{Medical Pipeline} est une application complète de \textbf{science des données médicales} qui couvre toute la chaîne de valeur d'un projet de Machine Learning~: de la \textbf{collecte automatisée de données} (web scraping du site officiel du NHS) jusqu'à un \textbf{outil de diagnostic interactif} présenté dans une interface web moderne.

À partir d'une liste de symptômes cochés par l'utilisateur, l'application prédit si un patient est \textbf{«~Malade~»} ou \textbf{«~Sain~»}, estime une \textbf{probabilité} et un \textbf{niveau de confiance}, puis identifie la \textbf{maladie la plus probable} par similarité symptomatique.

Le projet démontre la maîtrise d'un \textbf{pipeline de données de bout en bout}~:
\begin{center}
\textbf{Scraping $\rightarrow$ Nettoyage $\rightarrow$ Contrôle qualité $\rightarrow$ Entraînement $\rightarrow$ Prédiction $\rightarrow$ Interface.}
\end{center}

\begin{quote}
\textcolor{warning}{\textbf{Note importante~:}} l'outil est à vocation \textbf{éducative uniquement} et ne remplace en aucun cas un avis médical professionnel. Cet avertissement est intégré directement dans l'interface.
\end{quote}

% ================= 2. OBJECTIFS =================
\section{Objectifs du projet}

\begin{longtable}{@{}p{0.5cm} p{11.5cm} p{2cm}@{}}
\toprule
\textbf{\#} & \textbf{Objectif} & \textbf{Statut} \\
\midrule
\endhead
1 & Collecter automatiquement des données de symptômes depuis une source fiable (NHS) & \ok \\
2 & Construire un dataset structuré et exploitable (X = symptômes, y = label) & \ok \\
3 & Mettre en place un pipeline de nettoyage et de contrôle qualité & \ok \\
4 & Entraîner et comparer plusieurs modèles de classification & \ok \\
5 & Identifier la maladie probable au-delà du diagnostic binaire & \ok \\
6 & Proposer une interface utilisateur claire et accessible & \ok \\
7 & Assurer la reproductibilité et le versionnement (Git + DVC) & \partiel \\
\bottomrule
\end{longtable}

% ================= 3. ARCHITECTURE =================
\section{Architecture technique}

\subsection{Vue d'ensemble du pipeline}

\begin{center}
\fbox{\texttt{scraper\_nhs}} $\rightarrow$
\fbox{\texttt{cleaning}} $\rightarrow$
\fbox{\texttt{quality\_check}} $\rightarrow$
\fbox{\texttt{model}} $\rightarrow$
\fbox{\texttt{app}}
\end{center}

\begin{center}
\small (collecte) \quad (nettoyage) \quad (qualité) \quad (ML / IA) \quad (interface)\\
$\downarrow$ \texttt{data/raw/*.csv} \hspace{4cm} $\downarrow$ \texttt{models/*.pkl}
\end{center}

\subsection{Description des modules}

\begin{longtable}{@{}p{3cm} p{3cm} p{8cm}@{}}
\toprule
\textbf{Fichier} & \textbf{Rôle} & \textbf{Points clés} \\
\midrule
\endhead
\texttt{scraper\_nhs.py} & Collecte & Scrape les pages du NHS, détecte \textbf{50 symptômes cibles} via un dictionnaire de mots-clés, génère des profils «~sains~» synthétiques, exporte un CSV. \\
\texttt{cleaning.py} & Nettoyage & Suppression des doublons, gestion des valeurs manquantes (remplissage à 0), correction des types. \\
\texttt{quality\_check.py} & Qualité & 4 contrôles automatiques (manquants, doublons, équilibre des classes, complétude) avec un \textbf{score sur 100}. \\
\texttt{model.py} & ML / IA & Entraîne un \textbf{Random Forest} et une \textbf{Régression Logistique}, identifie la maladie par \textbf{similarité de Jaccard}. \\
\texttt{app.py} & Interface & Application \textbf{Streamlit} en 3 onglets, design sombre personnalisé en CSS. \\
\bottomrule
\end{longtable}

\subsection{Stack technologique}
\begin{itemize}
  \item \textbf{Langage}~: Python
  \item \textbf{Collecte}~: \texttt{requests}, \texttt{BeautifulSoup} (bs4)
  \item \textbf{Données}~: \texttt{pandas}, \texttt{numpy}
  \item \textbf{Machine Learning}~: \texttt{scikit-learn} (RandomForest, LogisticRegression, StandardScaler)
  \item \textbf{Interface}~: \texttt{Streamlit}
  \item \textbf{Sérialisation}~: \texttt{pickle}
  \item \textbf{Versionnement}~: \texttt{Git} + \texttt{DVC}
\end{itemize}

% ================= 4. FONCTIONNALITES =================
\section{Fonctionnalités réalisées}

\subsection{Collecte de données (Scraping NHS)}
\begin{itemize}
  \item Récupération automatique de la liste des maladies depuis \url{https://www.nhs.uk/conditions/}.
  \item Détection de \textbf{50 symptômes} par correspondance de mots-clés (ex. \textit{fever} $\leftrightarrow$ «~fever, high temperature, pyrexia~»).
  \item Extraction ciblée des sections «~Symptoms~» des pages.
  \item Génération de \textbf{profils sains synthétiques} pour équilibrer le dataset.
  \item Paramétrage du nombre de maladies (10 à 100) et export CSV automatique.
\end{itemize}

\subsection{Pipeline de traitement}
\begin{itemize}
  \item Nettoyage des doublons, valeurs manquantes et types.
  \item Contrôle qualité automatisé~: manquants $<$ 15\,\%, absence de doublons, équilibre malades/sains (30--70\,\%), au moins 1 symptôme par ligne.
\end{itemize}

\subsection{Modélisation et diagnostic}
\begin{itemize}
  \item Entraînement comparé de \textbf{2 modèles} avec \texttt{class\_weight="balanced"}.
  \item Séparation train/test (80/20) stratifiée.
  \item Prédiction avec \textbf{probabilité} et \textbf{niveau de confiance} (Haute / Moyenne / Faible).
  \item \textbf{Identification de la maladie probable} via similarité de Jaccard.
\end{itemize}

\subsection{Interface utilisateur}
\begin{itemize}
  \item Application web en \textbf{3 onglets} guidant l'utilisateur pas à pas.
  \item Visualisations~: métriques, top 10 des symptômes, cartes de résultats colorées.
  \item Design sombre soigné et avertissement médical.
\end{itemize}

% ================= 5. BILAN =================
\section{Bilan technique}

\subsection{Points forts}
\begin{itemize}
  \item \textbf{Pipeline complet et cohérent} couvrant tout le cycle de vie d'un projet ML.
  \item \textbf{Séparation claire des responsabilités} (un module = une étape).
  \item \textbf{Code lisible et documenté} (docstrings, nommage explicite).
  \item \textbf{Bonnes pratiques ML}~: stratification, gestion du déséquilibre, normalisation.
  \item \textbf{Valeur ajoutée} grâce à l'identification de maladie.
  \item \textbf{Interface professionnelle} et expérience utilisateur soignée.
\end{itemize}

\subsection{Limites et axes d'amélioration}

\begin{longtable}{@{}p{7cm} p{8cm}@{}}
\toprule
\textbf{Limite} & \textbf{Recommandation} \\
\midrule
\endhead
Absence de \texttt{requirements.txt} & Figer les dépendances pour la reproductibilité. \\
Pas de tests automatisés & Ajouter des tests unitaires (\texttt{pytest}). \\
Données «~saines~» synthétiques & Utiliser des données réelles ou une génération plus réaliste. \\
Détection par mots-clés simple & Explorer le NLP (embeddings) pour l'extraction. \\
Pas de README & Rédiger un \texttt{README.md} (installation, lancement). \\
Évaluation limitée à l'\textit{accuracy} & Ajouter précision, rappel, F1-score, matrice de confusion. \\
DVC non finalisé & Compléter le suivi des datasets/modèles. \\
Pas de CI/CD & Ajouter un workflow GitHub Actions (lint + tests). \\
\bottomrule
\end{longtable}

% ================= 6. INDICATEURS =================
\section{Indicateurs du projet}

\begin{longtable}{@{}p{9cm} p{5cm}@{}}
\toprule
\textbf{Indicateur} & \textbf{Valeur} \\
\midrule
\endhead
Nombre de modules Python & 5 \\
Symptômes pris en charge & 50 \\
Modèles de ML & 2 (Random Forest, Régression Logistique) \\
Modèles sérialisés (\texttt{models/}) & 5 fichiers \texttt{.pkl} \\
Contrôles qualité automatiques & 4 \\
Onglets de l'interface & 3 \\
Source de données & NHS (officielle) \\
Commits Git & 1 («~first commit~») \\
\bottomrule
\end{longtable}

% ================= 7. PERSPECTIVES =================
\section{Perspectives pour l'année prochaine}

\begin{enumerate}
  \item \textbf{Fiabiliser la base}~: \texttt{requirements.txt}, \texttt{README.md}, tests unitaires, CI/CD.
  \item \textbf{Enrichir les données}~: plus de maladies, données saines réelles, mise à jour automatisée.
  \item \textbf{Améliorer les modèles}~: métriques avancées, validation croisée, XGBoost, etc.
  \item \textbf{Renforcer l'extraction de symptômes} par des techniques de NLP.
  \item \textbf{Déploiement}~: héberger l'application (Streamlit Cloud, Docker).
  \item \textbf{Finaliser le versionnement} des données et modèles avec DVC.
\end{enumerate}

% ================= 8. CONCLUSION =================
\section{Conclusion}

Le projet \textbf{Medical Pipeline} constitue une \textbf{réalisation solide et aboutie} illustrant la mise en œuvre complète d'un projet de data science, depuis la collecte des données jusqu'à la livraison d'un produit utilisable. Il combine \textbf{compétences techniques} (scraping, ingénierie de données, machine learning) et \textbf{souci de l'expérience utilisateur}.

Les bases posées cette année sont saines et bien structurées. Les principaux chantiers de l'année à venir concernent la \textbf{professionnalisation} (tests, documentation, déploiement) et l'\textbf{amélioration de la qualité des données et des modèles}. Le projet est ainsi très bien positionné pour évoluer vers une version plus robuste et déployable.

\vspace{1.5cm}
\hrule
\vspace{0.3cm}
{\small\itshape Rapport généré le 10 juin 2026 -- Projet Medical Pipeline (ElKarizma). Outil à finalité éducative -- ne se substitue pas à un avis médical professionnel.}

\end{document}
